\documentclass[11pt,letterpaper]{article}

\usepackage[margin=1.9cm]{geometry}
\usepackage{fontspec}
\usepackage{graphicx}
\usepackage{xcolor}
\usepackage{textcomp}
\usepackage{amssymb}
\usepackage{tabularx}
\usepackage{array}
\usepackage{enumitem}
\usepackage{multicol}
\usepackage{titlesec}
\usepackage{fancyhdr}
\usepackage{lastpage}
\usepackage{ragged2e}
\usepackage{setspace}
\usepackage[most]{tcolorbox}
\usepackage[hidelinks]{hyperref}

\setmainfont{Arial}
\setsansfont{Arial}
\setmonofont{Consolas}

\setstretch{1.12}

% --- Paleta institucional ---
\definecolor{AIEPRed}{HTML}{D62027}
\definecolor{AIEPNavy}{HTML}{102A43}
\definecolor{AIEPInk}{HTML}{243B53}
\definecolor{AIEPSlate}{HTML}{52606D}
\definecolor{AIEPPaper}{HTML}{F8F3EC}
\definecolor{AIEPSoftBlue}{HTML}{E6EEF7}
\definecolor{AIEPGreen}{HTML}{3FAE6A}
\definecolor{AIEPGold}{HTML}{E0BC5A}
\definecolor{AIEPAmber}{HTML}{E8A33D}
\definecolor{Line}{HTML}{D8E1EA}

\pagestyle{fancy}
\fancyhf{}
\renewcommand{\headrulewidth}{0pt}
\fancyfoot[L]{\small\textcolor{AIEPSlate}{GeoGreen Escolar · Banco de ideas semilla}}
\fancyfoot[R]{\small\textcolor{AIEPSlate}{\thepage/\pageref{LastPage}}}

\setlength{\parindent}{0pt}
\setlength{\parskip}{6pt}
\setlist[itemize]{leftmargin=*,topsep=2pt,itemsep=3pt}
\setlist[enumerate]{leftmargin=*,topsep=2pt,itemsep=3pt}
\renewcommand{\arraystretch}{1.3}

\titleformat{\section}{\Large\bfseries\color{AIEPNavy}}{}{0pt}{}
\titlespacing{\section}{0pt}{16pt}{8pt}

\newcolumntype{Y}{>{\RaggedRight\arraybackslash}X}
\newcolumntype{L}[1]{>{\RaggedRight\arraybackslash}p{#1}}

\newtcolorbox{infobox}[2][]{
  enhanced, colback=#2!7, colframe=#2, boxrule=0.5pt, arc=1.4mm,
  left=4mm, right=4mm, top=3.5mm, bottom=3.5mm, #1
}

% --- Etiqueta del módulo (dato para quien arma el circuito) ---
\newcommand{\serig}[1]{%
  \tcbox[on line, colback=AIEPInk, colframe=AIEPInk, boxrule=0pt,
    arc=0.8mm, left=1.8mm, right=1.8mm, top=0.5mm, bottom=0.5mm]%
    {\small\bfseries\textcolor{white}{\texttt{#1}}}%
}

% --- Chip de nivel de dificultad ---
\newcommand{\nivel}[2]{%
  \tcbox[on line, colback=#1!14, colframe=#1, boxrule=0.6pt,
    arc=0.9mm, left=1.8mm, right=1.8mm, top=0.5mm, bottom=0.5mm]%
    {\footnotesize\bfseries\textcolor{#1}{#2}}%
}
\newcommand{\nbasico}{\nivel{AIEPGreen}{Fácil}}
\newcommand{\nmedio}{\nivel{AIEPAmber}{Intermedio}}
\newcommand{\nreto}{\nivel{AIEPRed}{Desafío}}

% --- Tarjeta de sensor (foto + nombre llano + etiqueta + uso) ---
% #1 archivo  #2 nombre llano  #3 etiqueta  #4 para qué sirve
\newcommand{\poolcard}[4]{%
\begin{tcolorbox}[colback=white, colframe=Line, boxrule=0.6pt, arc=1.4mm,
  left=1.5mm, right=1.5mm, top=2mm, bottom=2mm, halign=center, valign=top]
\includegraphics[height=2.2cm,keepaspectratio]{assets/#1}\\[1.4mm]
{\small\bfseries\color{AIEPNavy}#2}\\[0.6mm]
\serig{#3}\\[1.2mm]
{\footnotesize\textcolor{AIEPSlate}{#4}}
\end{tcolorbox}%
}

% --- Ficha de idea ---
% #1 número  #2 título  #3 nivel  #4 problema/para qué (ambiental)
% #5 sensores (entrada)  #6 cómo funciona  #7 variación o reto
\newtcolorbox{ideabox}{
  enhanced, colback=white, colframe=Line, boxrule=0.7pt, arc=1.6mm,
  left=4mm, right=4mm, top=3.5mm, bottom=3.5mm, breakable=false,
  before skip=9pt, after skip=9pt
}
\newcommand{\idea}[7]{%
\begin{ideabox}
{\large\bfseries\color{AIEPNavy}#1.\ #2}\hfill#3\par
\vspace{2.5mm}
{\color{AIEPSlate}\textbf{Para qué:}} #4\par\vspace{1.8mm}
{\color{AIEPSlate}\textbf{Qué usa:}} #5\par\vspace{1.8mm}
{\color{AIEPSlate}\textbf{Cómo funciona:}} #6\par\vspace{1.8mm}
{\color{AIEPGreen}\textbf{Sube de nivel:}} #7
\end{ideabox}%
}

\begin{document}

% ===================== PORTADA =====================
\begin{tcolorbox}[enhanced, colback=AIEPNavy, colframe=AIEPNavy, arc=0mm,
  boxrule=0pt, left=6mm, right=6mm, top=6mm, bottom=6mm]
{\color{white}\Huge\bfseries Banco de ideas semilla}\\[3mm]
{\color{white}\Large GeoGreen Escolar}\\[4mm]
{\color{AIEPSoftBlue}\large Proyectos de \textbf{sostenibilidad y medioambiente} listos para cuando
un equipo dice ``no se me ocurre nada''. Todos se arman con el kit que ya tenemos y siguen la misma
idea del prototipo GeoGreen: \textbf{el aparato mide algo, decide si pasó un límite y avisa con una
luz o una alarma}.}
\end{tcolorbox}

\begin{infobox}{AIEPSoftBlue}
\textbf{Cómo usar este banco.}
La idea \emph{no} es entregarle un proyecto cerrado al estudiante, sino \textbf{darle un empujón}.
Cuando un equipo no se decide: ofréceles \textbf{2 o 3} de estas ideas, pídeles que elijan una y
que la \textbf{adapten} (cambiar el lugar, el límite, el sensor o lo que hace). Así la propuesta
sigue siendo suya. Todas comparten el espíritu de GeoGreen: \textbf{cuidar el medioambiente y los
recursos}. Cada idea trae una pista de ``\textcolor{AIEPGreen}{sube de nivel}'' para los equipos
que quieran ir más allá. \emph{No hace falta saber programar para entender la idea}: cada ficha dice
qué mide y qué hace.
\end{infobox}

\begin{infobox}{AIEPGreen}
\textbf{El aparato primero; el software es el ``sube de nivel''.} El prototipo base es solo el
dispositivo físico (sensor + luz/alarma) y lo puede armar \emph{cualquier} equipo. Sumarle
software --- una app, un panel en vivo, un mapa o un aviso al celular --- es lo que lo lleva al
siguiente nivel, y conecta directo con la app de GeoGreen y la placa con WiFi. Por eso casi todas
las ideas tienen su capa de software en ``\textbf{sube de nivel}''.
\end{infobox}

% ===================== POOL DE SENSORES =====================
\section{Los sensores que vamos a usar}
El kit trae muchos sensores. Para empezar, estos son los más fáciles y los que más dan juego.
Un \textbf{sensor} es la parte que ``siente'' algo del mundo (distancia, luz, humedad\ldots); una
\textbf{salida} es la parte que ``avisa o actúa'' (una luz, un sonido, encender un aparato).
Con esta mano de piezas se arman las ideas de este banco (y muchas más que se les ocurran).
(La etiqueta oscura, tipo \serig{HC-SR04}, es solo el nombre técnico, útil para quien arma el
circuito.)

\vspace{3mm}
{\large\bfseries\color{AIEPNavy}Sensores (lo que el proyecto ``siente'')}
\vspace{2mm}

\begin{tcbraster}[raster columns=4, raster equal height=rows,
  raster column skip=2.6mm, raster row skip=2.6mm]
\poolcard{hc-sr04.jpg}{Distancia}{HC-SR04}{Mide qué tan lejos está algo}
\poolcard{ky-015.jpg}{Temperatura y humedad}{KY-015}{Del aire de un lugar}
\poolcard{ky-001.jpg}{Temperatura}{KY-001}{Precisa, incluso del agua}
\poolcard{ky-018.jpg}{Luz}{KY-018}{Cuánta luz hay (día/noche)}
\poolcard{soil-moisture.jpg}{Humedad de la tierra}{Soil}{Si la tierra está seca}
\poolcard{water-level.jpg}{Nivel de agua}{Water}{Cuánta agua hay}
\poolcard{ky-021.jpg}{Sensor magnético}{KY-021}{Si una tapa quedó abierta}
\end{tcbraster}

\vspace{3mm}
{\large\bfseries\color{AIEPGreen}Salidas (cómo el proyecto ``avisa o actúa'')}\\
{\small\textcolor{AIEPSlate}{Las mismas tres del prototipo GeoGreen.}}
\vspace{2mm}

\begin{tcbraster}[raster columns=4, raster equal height=rows,
  raster column skip=2.6mm, raster row skip=2.6mm]
\poolcard{ky-011.jpg}{Luz de colores}{KY-011}{Semáforo verde / rojo}
\poolcard{ky-012.jpg}{Zumbador}{KY-012}{Alarma sonora}
\poolcard{ky-019.jpg}{Relé}{KY-019}{Enciende un aparato real}
\end{tcbraster}

\begin{infobox}{AIEPSoftBlue}
Existe además un \textbf{manual aparte con la foto y los datos de cada sensor} del kit, por si un
equipo quiere usar uno que no está en esta lista.
\end{infobox}

\begin{infobox}{AIEPRed}
\textbf{Seguridad.} El sensor de distancia funciona a 5\,V y se conecta directo a nuestra placa
(Arduino UNO R4). El relé puede manejar aparatos de 220\,V: si se llega a ese punto, siempre
con \textbf{supervisión docente}.
\end{infobox}

% ===================== IDEAS =====================
\section{12 ideas con impacto territorial y social}
El sello de AIEP es el \textbf{impacto territorial y social}. Por eso cada idea parte de un
\textbf{problema real de Osorno y la región de Los Lagos} y crece hacia un mapa o una red que llega
a toda la comunidad. El aparato es simple; lo que escala es el dato compartido.

\begin{infobox}{AIEPGold}
\textbf{De la maqueta al territorio.} Cada idea arranca como un prototipo \textbf{simple y barato}
--- un sensor, un \textbf{umbral} y una alarma --- pero su gracia es que \textbf{escala}: cuando
muchos comparten el dato en un mapa o panel, el mismo aparato chico se vuelve un sistema que sirve a
un barrio, una ciudad o toda la región. Empieza pequeño, piensa en grande.
\end{infobox}

\subsection*{\large\textcolor{AIEPGreen}{Aire y calefacción (la crisis de Osorno)}}

\idea{1}{¿Está seca tu leña?}{\nbasico}%
{Osorno es una de las ciudades con peor aire de Chile por la leña húmeda: al quemarse, llena el
aire de humo dañino. Saber si tu leña está lista para quemar limpio baja el humo y cuida la salud
de todo el barrio.}%
{\serig{Soil} clavado en un trozo de leña mide qué tan húmeda está.}%
{Compara tu leña con una muestra bien seca: si se acerca a la seca, verde, lista para quemar
limpio; si sigue húmeda, rojo, hay que secarla más.}%
{Un mapa del barrio con puntos donde se consigue leña seca y campañas para bajar el humo entre todos.}

\idea{2}{Detector de humedad y moho en casa}{\nbasico}%
{Las viviendas frías y húmedas del sur enferman (asma, bronquitis) y crían moho. Avisar a tiempo
cuándo ventilar es salud para miles de familias.}%
{\serig{KY-015} mide la humedad del aire de la pieza.}%
{Si la humedad se mantiene sobre $\sim$70\% por más de media hora: rojo, ``ventila''. Verde cuando
baja a un nivel sano.}%
{Graficar la humedad del día y armar un ranking de las salas o viviendas más húmedas de un sector.}

\subsection*{\large\textcolor{AIEPGreen}{Agua: crisis hídrica y comunidades}}

\idea{3}{Vigía del estanque de Agua Potable Rural}{\nmedio}%
{En la región, cientos de sistemas de Agua Potable Rural apenas alcanzan a abastecer a sus
comunidades; quedarse sin agua sin aviso es grave. Vigilar el nivel del estanque permite reaccionar
a tiempo.}%
{\serig{HC-SR04} montado arriba del estanque, midiendo la distancia al agua.}%
{El nivel sale de la distancia y la altura del estanque (lleno = poca distancia). Verde sobre 60\%,
amarillo 30--60\%, rojo bajo 30\% + alarma.}%
{Avisar al operador por celular y mostrar el nivel de varios estanques de la zona en un panel para
coordinar el reparto.}

\idea{4}{Riego eficiente para la sequía}{\nbasico}%
{La región fue declarada en emergencia agrícola por falta de agua. Regar solo cuando la tierra lo
necesita cuida un recurso escaso; ideal para huertos escolares o invernaderos comunitarios.}%
{\serig{Soil} mide la humedad de la tierra del cultivo o el huerto.}%
{Bajo el umbral de ``tierra seca'': aviso de regar (o riega solo con el relé \serig{KY-019}).
Húmeda: verde, no riega.}%
{Un panel del huerto con varias zonas y aviso al celular de cuál necesita agua, para repartir el
riego sin desperdicio.}

\idea{5}{Alarma de reserva de agua}{\nbasico}%
{En una escuela rural o una sede, quedarse sin agua en el estanque o los bidones detiene todo. Una
alarma simple de ``queda poca'' evita el apuro.}%
{\serig{Water} puesto a la altura del nivel mínimo del estanque o bidón.}%
{Mientras toca agua: verde. Cuando el agua baja y deja de tocarlo: rojo + alarma ``rellenar''.}%
{Avisar por celular y, con un segundo punto, distinguir ``lleno / a la mitad / vacío''.}

\subsection*{\large\textcolor{AIEPGreen}{Ríos y clima}}

\idea{6}{Alerta temprana de crecida de río}{\nmedio}%
{Los ríos del sur (como el Rahue y el Damas) se desbordan con las lluvias. Un aviso temprano de que
el agua sube da tiempo a la comunidad para reaccionar.}%
{\serig{HC-SR04} instalado en un puente, mirando hacia abajo, hacia el agua.}%
{Si el agua sube, la distancia baja. Verde normal, amarillo alerta, rojo + alarma en nivel
peligroso, a tiempo para avisar a la comunidad.}%
{Enviar el nivel al celular o a una web en vivo y mostrar varios puntos del cauce en un mapa de
alerta comunitaria.}

\idea{7}{Salud del río}{\nmedio}%
{El Río Damas y otros sufren descargas que dañan la vida del agua. Vigilar el río y delatar cuándo
algo cambia ayuda a defenderlo y a exigir que se cuide.}%
{\serig{KY-001} sumergido mide la temperatura del agua.}%
{Compara con la temperatura normal del río; una subida brusca delata una descarga río arriba y
enciende la alerta.}%
{Medir en varios puntos del río y armar un mapa ciudadano de la salud del agua.}

\subsection*{\large\textcolor{AIEPGreen}{Residuos}}

\idea{8}{Contenedor inteligente (el caso base)}{\nbasico}%
{Saber cuándo un basurero está lleno sin abrirlo: menos viajes de camión y menos desbordes en la
ciudad.}%
{\serig{HC-SR04} mide la distancia de la tapa a la basura.}%
{La distancia se vuelve \% de llenado. Verde bajo 40\%, amarillo 40--80\%, rojo sobre 80\% +
alarma. Es el prototipo de GeoGreen.}%
{Mostrar todos los contenedores en un mapa de la ciudad (la app de GeoGreen) y optimizar la ruta de
recolección.}

\idea{9}{Alerta de tapa abierta en contenedores}{\nbasico}%
{Un contenedor o punto limpio con la tapa abierta junta lluvia, da malos olores y atrae plagas.
Avisar para cerrarlo cuida el barrio.}%
{\serig{KY-021} (sensor magnético) más un imán pegado a la tapa.}%
{Tapa cerrada (imán cerca): verde. Abierta más de 2 minutos: amarillo; más de 10: rojo + alarma.}%
{Un mapa de los contenedores que quedan abiertos seguido, para enfocar el mantenimiento o una campaña.}

\subsection*{\large\textcolor{AIEPGreen}{Energía e infraestructura}}

\idea{10}{Alumbrado público que avisa cuándo falla}{\nmedio}%
{Una ciudad tiene miles de luminarias; muchas quedan quemadas de noche o encendidas de día sin que
nadie lo note. Detectarlo ahorra energía pública y mantiene las calles seguras.}%
{\serig{KY-018} junto a la luminaria.}%
{De noche con poca luz y la luminaria apagada: falla (quemada). De día con mucha luz y la luminaria
encendida: desperdicio. En ambos casos, rojo.}%
{Un mapa de la ciudad con cada luminaria en verde o rojo, para que mantenimiento sepa al instante
cuál arreglar.}

\idea{11}{¿Dónde pega más el sol?}{\nmedio}%
{Aprovechar el sol --- para un huerto, un panel solar o secar leña --- parte por saber qué lugar
recibe más luz a lo largo del día. Una medición simple lo resuelve y evita decisiones a ojo.}%
{\serig{KY-018} mide la luz en distintos puntos del patio o el techo.}%
{Compara cuánta luz recibe cada lugar y marca el más soleado; también delata salas que gastan luz
de día por falta de luz natural.}%
{Registrar el sol por hora y dibujar un ``mapa de sol'' del colegio o el terreno.}

\subsection*{\large\textcolor{AIEPGreen}{Producción local}}

\idea{12}{Guardián de la cadena de frío (leche, vacunas, alimentos)}{\nbasico}%
{En una región lechera y rural, mantener el frío de la leche, las vacunas o los alimentos es clave:
si el frío se corta, se pierde un producto caro y se arriesga la salud.}%
{\serig{KY-001} dentro del refrigerador, el tanque de leche o el cooler.}%
{Cada caso tiene su rango (ej. la leche cruda bajo $\sim$4\,\textdegree C). Si la temperatura sale
del rango seguro: rojo + alarma al instante.}%
{Guardar el historial de temperatura y avisar al teléfono del encargado apenas algo falla.}

% ===================== FICHA TEMPLATE =====================
\section{Ficha de idea (resumen en 4 líneas)}
Cualquier proyecto, propio o tomado de este banco, se puede describir en 4 líneas. Si el equipo
llena esta ficha, ya tiene su idea tecnológica inicial:

\begin{infobox}{AIEPGold}
\setstretch{1.5}
\textbf{1. Problema ambiental / para qué sirve:} \dotfill\par
\textbf{2. Qué mide y con qué sensor:} \dotfill\par
\textbf{3. Cuándo se activa (el límite):} \dotfill\ (verde / amarillo / rojo)\par
\textbf{4. Cómo avisa o actúa:} \dotfill\ (luz, alarma, encender algo, contar\ldots)
\end{infobox}

\begin{infobox}{AIEPSoftBlue}
\textbf{Regla de oro para destrabar.} Casi todo proyecto entra en una frase:
``\emph{cuando} [el sensor] pase de [un límite], \emph{entonces} [avisa o acciona]''.
Si el equipo completa esa frase, ya tiene proyecto.
\end{infobox}

\end{document}
